\section{Integer intervals and the two growth regimes}
\label{sec:transfers}

The cyclic coloring has a deterministic product extension.  The relevant
digit is the first place at which the common difference is nonzero;
division by the corresponding power of the base then introduces no carry.
The related digit construction of Erd\H{o}s and Tur\'an
\cite[p.~263]{erdosTuran1936} is recalled, with the least-nonzero-digit
argument, by Fox and Hunter~\cite[Section~2]{foxHunter2026}.
We prove the arbitrary-base product needed here directly.

\begin{proposition}[Digit product]
\label{transfer:product}
Let $N,k\ge2$ be integers.  Suppose that a coloring
$C:\Z/N\Z\to\{0,1\}$ has no monochromatic progression
$a,a+d,\ldots,a+(k-1)d$ with $d\ne0$ in $\Z/N\Z$.
Then, for every positive integer $m$,
\[
 W_{2^m}(k)>N^m.
\]
Consequently $W_r(k)>N^{\lfloor\log_2r\rfloor}$ for every integer
$r\ge2$.
\end{proposition}

\begin{proof}
For $0\le n<N^m$, write its base-$N$ expansion as
\[
 n=\sum_{i=0}^{m-1}n_iN^i,\qquad 0\le n_i<N,
\]
and assign the color
\[
 \bigl(C(n_0),\ldots,C(n_{m-1})\bigr)\in\{0,1\}^m.
\]
Here and below a digit is identified with its residue modulo the base.
Suppose that $a,a+d,\ldots,a+(k-1)d$ is a monochromatic progression
in this interval, with $d>0$.  Since $k\ge2$, we have $d<N^m$.
There is therefore a largest integer $i\in\{0,\ldots,m-1\}$ for which
$N^i\mid d$.  Put $e=d/N^i$, so that $N\nmid e$.
For every $j\in\{0,\ldots,k-1\}$, the exact identity
\begin{equation}
 \left\lfloor\frac{a+jd}{N^i}\right\rfloor
 =\left\lfloor\frac a{N^i}\right\rfloor+je
 \label{transfer:division}
\end{equation}
holds because $je$ is an integer.  Reducing
Equation~\eqref{transfer:division} modulo $N$ shows that the $i$th
digits form a cyclic progression with nonzero step $e\pmod N$.
Monochromaticity of the color tuples would make $C$ constant on that
progression, a contradiction.

Translation by one now gives a coloring of $[N^m]$ with $2^m$ colors
and no monochromatic $k$-term progression.  This proves the first
assertion.  For the second, take $m=\lfloor\log_2r\rfloor\ge1$ and
embed the at most $2^m\le r$ color tuples into $[r]$.
\end{proof}

In particular, $m=1$ restricts the cyclic coloring to an ordinary
interval of length $N$.  The product uses the same fixed coloring $C$
in every digit, for arbitrarily many digits.

\begin{proof}[Proof of Theorem~\ref{main:intro}]
Theorem~\ref{perturb:cyclic} supplies, for every integer $k$ above an
absolute threshold, such a cyclic coloring with
\[
 N=q^D\ge k^{ck},\qquad c=10^{-5}.
\]
Proposition~\ref{transfer:product} gives, simultaneously for every
integer $r\ge2$,
\[
 W_r(k)>N^{\lfloor\log_2r\rfloor}
 \ge k^{ck\lfloor\log_2r\rfloor}.
\]
The threshold on $k$ was fixed before $r$ was introduced, so it is
absolute and uniform in the number of colors.
\end{proof}

\subsection{Growth in the number of colors}

For fixed progression length, a separate construction gives a useful
bound with a threshold independent of the length.  The construction
records a coarse interval for each digit, to force midpoint relations
to hold coordinate by coordinate, and a squared norm, to exclude
nonconstant midpoints. This is the digit-and-sphere geometry of
Behrend~\cite{behrend1946}, used here to partition the entire interval
into progression-free color classes.

\begin{theorem}
\label{transfer:large-r}
For every pair of integers $r\ge256$ and $k\ge3$,
\[
 W_r(k)>\exp\!\left(\frac{(\log r)^2}{64\log2}\right).
\]
\end{theorem}

\begin{proof}
Set
\[
 b=\lfloor r^{1/4}\rfloor,\qquad
 s=\left\lfloor\frac{\log r}{4\log2}\right\rfloor.
\]
Both are positive, and $b\ge2$.  Partition the digit set
$\{0,\ldots,b-1\}$ into
\[
 H_0=\{0,\ldots,\lceil b/2\rceil-1\},\qquad
 H_1=\{\lceil b/2\rceil,\ldots,b-1\},
\]
and let $h(u)\in\{0,1\}$ record the half containing $u$.
An integer $0\le n<b^s$ has a unique digit vector
$x=(x_0,\ldots,x_{s-1})$ with
$n=\sum_{i=0}^{s-1}x_i b^i$.  Give it the color
\begin{equation}
 \left(h(x_0),\ldots,h(x_{s-1}),\ \sum_{i=0}^{s-1}x_i^2\right).
 \label{transfer:norm-color}
\end{equation}
There are at most $2^s(1+s(b-1)^2)\le2^s(1+sb^2)$ such colors.
For every integer $s\ge1$ we have $s+1\le2^s$, while
$2^s\le r^{1/4}$ and $b^2\le r^{1/2}$.  Hence
\begin{equation}
 2^s(1+sb^2)
 \le r^{1/4}(s+1)r^{1/2}
 \le r.
 \label{transfer:color-count}
\end{equation}

We show that this coloring contains no nonconstant monochromatic
three-term progression.  Suppose that three equally spaced integers
have digit vectors $x,y,z$ and the same color.  Their midpoint relation
is
\begin{equation}
 \sum_{i=0}^{s-1}(x_i-2y_i+z_i)b^i=0.
 \label{transfer:midpoint}
\end{equation}
For each $i$, the three digits $x_i,y_i,z_i$ lie in one half $H_0$ or
$H_1$.  Each half has diameter at most $\lceil b/2\rceil-1$, so
\[
 |x_i-2y_i+z_i|
 \le2(\lceil b/2\rceil-1)
 =\begin{cases}
 b-2,&\text{if $b$ is even},\\
 b-1,&\text{if $b$ is odd}.
 \end{cases}
\]
In particular this absolute value is strictly less than $b$.
Reducing Equation~\eqref{transfer:midpoint} modulo $b$ forces its
zeroth coefficient to be zero: it is divisible by $b$ and has absolute
value less than $b$.  Dividing by $b$ and repeating proves
$x_i-2y_i+z_i=0$ for every $i$.  Thus $x+z=2y$ in $\R^s$.
Equality of the final entries in Equation~\eqref{transfer:norm-color}
also gives $\norm{x}_2^2=\norm{y}_2^2=\norm{z}_2^2$, and therefore
\[
 \norm{x-z}_2^2
 =2\norm{x}_2^2+2\norm{z}_2^2-4\norm{y}_2^2=0.
\]
It follows that $x=y=z$, so the original three integers were equal.
This excludes every progression with positive difference and length
$k\ge3$, since its first three terms would already give a forbidden
triple.  After translating the interval by one and using
Equation~\eqref{transfer:color-count}, we obtain $W_r(k)>b^s$.

Finally let $L=\log r\ge8\log2$.  Since
$\lfloor u\rfloor\ge u/2$ for $u\ge1$, the parameter choices give
\[
 s\ge\frac{L}{8\log2},\qquad
 \log b\ge\frac L4-\log2\ge\frac L8.
\]
Consequently $\log(b^s)\ge L^2/(64\log2)$, proving the stated bound
for every $r\ge256$, independently of $k\ge3$.
\end{proof}

\subsection{Uniform growth statements}

\begin{corollary}
\label{transfer:limits}
With integer parameters in both infima,
\[
 \lim_{k\to\infty}
 \inf_{\substack{r\in\Z\\r\ge2}}
 \frac{\log W_r(k)}{k\log r}=\infty,
 \qquad
 \lim_{r\to\infty}
 \inf_{\substack{k\in\Z\\k\ge3}}
 \frac{\log W_r(k)}{\log r}=\infty.
\]
In particular, $W_r(k)^{1/k}\to\infty$ for every fixed integer
$r\ge2$, and $W_r(k)/r^A\to\infty$ as $r\to\infty$ for every fixed
integer $k\ge3$ and every fixed real $A>0$.
\end{corollary}

\begin{proof}
For $r\ge2$, the inequality
$\lfloor\log_2r\rfloor\ge(\log r)/(2\log2)$ and
Theorem~\ref{main:intro} give, for all $k\ge K_0$,
\[
 \inf_{\substack{r\in\Z\\r\ge2}}
 \frac{\log W_r(k)}{k\log r}
 \ge\frac{c}{2\log2}\log k.
\]
Theorem~\ref{transfer:large-r} gives, for all $r\ge256$,
\[
 \inf_{\substack{k\in\Z\\k\ge3}}
 \frac{\log W_r(k)}{\log r}
 \ge\frac{\log r}{64\log2}.
\]
The right-hand sides tend to infinity.  For fixed $r\ge2$,
Theorem~\ref{main:intro} also gives
$W_r(k)^{1/k}>k^{c\lfloor\log_2r\rfloor}\to\infty$.
For fixed $A>0$, Theorem~\ref{transfer:large-r} gives
\[
 \log\frac{W_r(k)}{r^A}
 >\frac{(\log r)^2}{64\log2}-A\log r\longrightarrow\infty,
\]
uniformly for integers $k\ge3$.
\end{proof}

The restrictions on the parameters are necessary: the exact values
$W_1(k)=k$, $W_r(1)=1$, and $W_r(2)=r+1$ are established along with
finiteness in Appendix~\ref{sec:upper}.
